\section*{Заключение}
\addcontentsline{toc}{section}{Заключение}
\label{sec:conclusion}
\draft{Написать после обсуждения результатов: подтверждённая точность, выигрыш скорости и достигнутое согласование наблюдений; затем ограничения восстановления истории потока и направления дальнейшей работы.}

Сформулирована модель двумерного теплопереноса в осадочном бассейне с заданной историей геометрии. В модели учитываются зависимость свойств от температуры, линейное изменение пористости по мощности слоя и времени, постепенное осадконакопление и разломные смещения. Начальное температурное поле определяется стационарной задачей с начальными граничными условиями. Отражательная способность витринита вычисляется по температурной истории материала от момента его отложения. Обратная задача состоит в определении пространственно постоянного базального потока при ограничениях на его монотонность и диапазон.

Разработан численный алгоритм на основе метода конечных объёмов на подвижной колоночной сетке. Тепловые потоки через вертикальные подграни определяются многоточечной аппроксимацией, учитывающей несовпадение уровней соседних ячеек. Для интегрирования по времени применяется неявный метод Эйлера, для решения нелинейных и линейных систем~--- методы Пикара и GMRES(100) с ILU(2). В обратной задаче все 17~параметров потока подбираются методом POUNDERS без вычисления градиента. Каждое новое вычисление вектора невязок требует полного прямого расчёта, включая начальное стационарное поле.

Программная реализация обеспечивает подготовку исходных данных, решение обеих задач и вывод полей на исходной геометрии. Однократная подготовка геометрии и сохранение рабочих матриц и векторов сокращают затраты при повторных прямых расчётах. На рассмотренных моделях при одном вычислительном потоке GMRES/ILU(2) уменьшает время линейного решения в 1,44--4,05~раза, полного прямого расчёта~--- в 1,31--2,55~раза относительно UMFPACK. Табулирование ускоряет вычисление аррениусовских интегралов в 6,55~раза. В проверке на аналитической температурной истории дополнительное относительное различие $R_o$ по сравнению с непосредственным вычислением первообразной не превышает $1{,}4\cdot10^{-8}\%$. Упрощение геометрии Модели~2 с допуском 5~м сокращает число ячеек с 120335 до 3943 при относительных различиях средних по слоям температуры и отражательной способности 1,31 и 1,12\% относительно расчёта без упрощения.

На гладких стационарных изготовленных решениях с постоянной и переменной теплопроводностью получена сходимость, близкая ко второму порядку по пространству. При совместном уменьшении пространственных шагов вдвое и временного шага вчетверо ошибка снижается в 3,99--4,00~раза как на неподвижном, так и на движущемся слое. Для наклонного контакта с разрывом теплопроводности наблюдается порядок около первого. Дополнительные проверки локализуют существенную ошибку потока на вертикальных подгранях между разными материалами. Тем самым установлены условия, в которых получена оценка второго порядка, и выявлено ограничение применяемой аппроксимации контактов.

В синтетической обратной задаче максимальные относительные невязки наблюдений составляют 0,0029\% по температуре и 0,0391\% по отражательной способности, а современные поля близки во всей области. Вместе с тем различие раннего потока достигает 11,01\%, и поля в прошлых конфигурациях различаются сильнее, чем в современности. Зависимость этих различий от возраста немонотонна. Полученный результат указывает на слабую различимость части параметров истории потока по современным наблюдениям.

Для Модели~1 относительные нормы различий средних по слоям температуры и отражательной способности с эталонным решением за все конфигурации равны 2,803 и 1,859\% соответственно. В обратной задаче Модели~3 за 10~итераций и 44~прямых расчёта найден поток, при котором в последующем контрольном расчёте согласуются в пределах 2\% 456 из 500~температур и 39 из 50~значений отражательной способности. Среднеквадратичные относительные невязки составляют 1,283 и 2,109\%. Таким образом, найденное решение улучшает согласование данных, но не обеспечивает попадания всех наблюдений в интервалы погрешности.

Дальнейшего исследования требуют аппроксимация потоков на наклонных контактах и устойчивость восстановления многопараметрической истории потока. Проверенные сгущения сетки и временного шага не устраняют оставшиеся расхождения со скважинными наблюдениями. Следующими задачами являются уточнение межслойных потоков, проверка более компактной параметризации, влияния начального приближения и весов данных. При оценке результатов необходимо по-прежнему различать точность прямого расчёта, согласование наблюдений и точность восстановленных параметров.
